% 22. Рекомендации
\newcommand{\recrow}[3]{%
  \begin{wcard}[top=2.4mm,bottom=2.4mm]
    \iconitem{navy}{#1}{{\bfseries\small\color{navy}#2}\par\vspace{0.6mm}\desc{#3}}
  \end{wcard}\par\vspace{3mm}}
\begin{frame}{Рекомендации по защите}
\framesubtitle{Что авторы советуют разработчикам детекторов}
\vspace{1.5mm}
\recrow{ic_shieldv_w.png}{Traffic-space adversarial training}{обучать на пакетных возмущениях, а не на $L_p$-норме признаков}
\recrow{ic_layers_w.png}{Сертификация в пространстве трафика}{ограничивать число вставленных пакетов, а не норму}
\recrow{ic_random_w.png}{Случайность на инференсе}{менять архитектуру/параметры $\rightarrow$ атакующий бьёт по «распределению моделей»}
\vspace{0.5mm}
\begin{darkcard}[halign=flush center,top=2mm,bottom=2mm]
  \bfseries Общий вывод: устойчивость ML-детекторов трафика остаётся нерешённой проблемой.
\end{darkcard}
\end{frame}
